%% The manuscript body, shared by every binding of this paper.
%% The class, the front matter and the bibliography style belong to the
%% binding; everything a reader reads belongs here, once.
\section{Introduction}

A small evaluation reports that one condition beats another. Underneath that
sentence is a table of questions, each marked right or wrong, and the sentence
is true because of how a few of those marks fell. How few is a question with an
exact answer, and it is almost never asked.

This paper asks it of one completed evaluation. A question set of \Questions{}
items was answered by an instruction model under \ConditionsWord{} retrieval
conditions and graded against gold strings, twice over --- once on widely
documented entities and once on less documented ones --- for \Answers{} graded
answers. The project that produced those runs computed \Contrasts{} paired
contrasts from them and recorded which it considered to have found something.
The answers are used exactly as recorded. What this paper adds to them is a set
of counts, identities, and a simulation: for each conclusion, the smallest
number of individual gradings that would have to have come out the other way
for it to stop holding; a closed form that matches the search on every
four-cell state of the three designs; the same three nested string rules
counted on all \Answers{} answers rather than on the image subset; and the
operating characteristic of that count under a paired Bernoulli model.

The count is exact rather than estimated. A paired contrast between two
conditions reduces to four numbers --- how many questions both got right, how
many each got right alone, how many neither got right --- and every quantity
computed from it, including the exact test, is a function of those four. One
grading change moves one question from one of the four to an adjacent one. So
the tables reachable by $k$ grading changes are exactly those reachable in $k$
steps, the space is small, and a two-index closed form returns the true minimum.
A breadth-first walk over the same lattice agrees with that form on every
four-cell state of the three designs (\ClosedFormStates{} states,
\ClosedFormMismatches{} mismatches) and is retained as the build's unit test.

Three conclusions were named before the search was written, and the answers are
not alike. The reported finding that the multimodal condition trails the
text-only one on image-dependent questions rests on \HarmK{} grading, and
\HarmCandidates{} of the \HarmGradings{} individual gradings on that subset
would each on their own erase it. The reported \HelpDiff{}-point difference in
the other direction has exact \(p=\HelpP{}\) and is reported as not significant;
it needs \HelpK{} gradings to reverse. Under the stored grader,
\SurvivesHolm{} of \Contrasts{} contrasts survive Holm correction, two of them
identical question by question; the long-tail text pair that does so needs
\SigK{}. The prespecified kill for this study was that all three come out at
\KillThreshold{} or more, which would have refuted the premise that these
conclusions turn on individual records; it \KillState{}, and the smallest is
\SmallestK{}.

Two things make this more than an arithmetic exercise. The first is that the
single grading is a disagreement already present in the record rather than a
hypothetical one. An independent regrade of the same answers by three nested
string rules finds \RegradeContested{} answers they disagree about. Of
those \RegradeContested{}, \RegradeImageFlip{} is among the \HarmCandidates{} that carry the harm
contrast: applied as the record states it, the reported difference moves from
\HarmDiff{} to \RegradedDiff{} percentage points. The conclusion and its
negation are separated by one adjudication of one string. The other three sit
on year-valued long-tail answers and, under the numeric-tolerant rule, are
enough to drop the Holm-surviving long-tail text pair out of the family of
\Contrasts{}.

The second is that a contrast carries two conclusions, not one, and they are not
equally secure. The \HelpDiff{}-point difference (exact \(p=\HelpP{}\), reported
as not significant) needs \HelpK{} gradings to change direction, which is
comparatively sturdy at this scale. Its verdict --- that an exact test does
not call it significant --- needs \HelpVerdictK{}. On the same four counts,
from the same \SubsetN{} questions, the direction is \HelpRatio{} times further
from being overturned than the significance statement attached to it. Among the
\Directional{} contrasts that have a direction to reverse at all, the two
distances differ by \DecoupledGap{} gradings or more in \DecoupledCount{} of
them. Reporting a contrast as ``not significant'' and reporting it as
``smaller'' are two different claims with two different supports, and nothing
in a conventional results table distinguishes them.

A third observation falls out of the same rows. Of the \Contrasts{} contrasts,
\Incapable{} have a smallest attainable p-value above the threshold: the two
conditions disagreed on so few questions that no assignment of those
disagreements could have produced a significant result. Their non-significance
is a fact about the question set and carries no information about the
conditions being compared. A further \AtFloor{} returned exactly the most
extreme value available to them, so they were at the limit of what the design
could say.

The object throughout is the reporting of the evaluation, and the comparison
the evaluation was built to make is left where its records leave it. The
reference benchmark this line of work would be measured against was never
obtained; that comparison is recorded as \PrimaryTier{}, the comparison against
the strongest published method that would follow it as \SotaTier{}, and the
counts below are read beside those states rather than in place of them. Whether
the retrieval conditions differ is a question these records do not settle, and
the two contrasts examined most closely enter as objects whose fragility is
measured.

The contribution has three parts. The first is a statistic: the fragility index
of the clinical-trial literature, carried to a paired grading unit and measured
against both claims a contrast carries, its direction and its verdict, with a
two-index closed form that exhaustive search confirms on every reachable table.
The second is that statistic's operating characteristic at small \(n\),
tabulated under a paired Bernoulli model. The third is an audit of one completed
evaluation with it, in which the survival of a Holm-corrected result is read
conditional on the grading rule and on the family it was corrected in. The
construction itself is Feinstein's; what is added is the unit it is counted in,
the second claim it is measured against, the closed form, and the operating
characteristic. At this scale a conclusion is a property of individual
gradings, and a results table that prints a difference and a p-value without
the number of gradings each rests on has left out the part a reader needs in
order to know what they are looking at.

\section{Related work}

\subsection{Counting the observations a significant result rests on}

The idea that a verdict should be reported alongside the number of observations
that sustain it originates in clinical epidemiology. Feinstein proposed
appraising a contrast of two proportions by asking how many patients would have
to be reclassified before the contrast stopped being significant, and named the
answer a unit fragility index~\cite{feinstein1990unit}. Walsh and colleagues
turned the idea into a survey instrument: across \(399\) randomised controlled
trials reporting at least one significant dichotomous outcome, they moved
patients in the smaller event group from non-event to event until Fisher's exact
test crossed \(0.05\), and found a median of eight~\cite{walsh2014statistical}.
A quarter of those trials required three or fewer. The statistic has been
criticised as a monotone re-expression of the p-value that penalises small
trials for being small~\cite{carter2017fragility}, and that criticism is worth
stating precisely, because it delimits what any counting statistic can claim: a
count is not an independent line of evidence about the effect, it is a
restatement of the same evidence on a scale a reader can hold in mind.
Machine learning has since reused the same name for a different object. Yang,
Zhang, Cao and colleagues define a Fragility Index that penalises large-error
binary predictions rather than counting how many dichotomous outcomes would
have to change~\cite{yang2023fragility}. That construction is a risk metric
rather than a perturbation count; the two share a name and nothing else, and
the statistic applied here is Feinstein's.

The meta-analytic ancestor is older still. Rosenthal asked how many unpublished
null studies would have to be sitting in file drawers before a significant
pooled result became non-significant, and reported that tolerance as a
number~\cite{rosenthal1979file}. The shared move in all three is to convert a
threshold crossing into a count of the units that would have to differ, and the
shared motivation is that a count communicates robustness in a way a p-value
near a threshold does not.

Three features distinguish the setting treated here from all of them. Their
designs are unpaired: two arms, two independent groups of units. Their search is
one-directional and incremental, walking the count upward until the threshold is
crossed, which returns a value that is attained but not certified as smallest.
And their question is exclusively about the verdict, because in a trial the
direction of an effect is not what the reader is being asked to trust. A paired
evaluation inverts the last of these: the direction of the difference is
frequently the sentence a reader carries away, and it can be fragile
independently of the verdict attached to it.

\subsection{Significance testing in retrieval and language evaluation}

The retrieval community has examined the reliability of its own comparisons for
decades. Zobel asked how far large-scale experimental results
generalise~\cite{zobel1998reliable}, and Voorhees showed that although
individual relevance judgments vary widely between assessors, system
\emph{rankings} built from different judgment sets correlate
highly~\cite{voorhees2000variations}. That finding is important context for what
follows, and it cuts in a specific direction: stability of a ranking over many
topics is a different property from stability of a single contrast over a few
dozen gradings, and the first does not imply the second.

Which test to apply, and what it licenses, has its own literature. Smucker,
Allan and Carterette compared the common choices and argued for randomisation
and bootstrap over the t-test's parametric
competitors~\cite{smucker2007comparison}; Urbano, Marrero and Martín examined
their optimality~\cite{urbano2013comparison}. Carterette treated multiplicity as
the central threat when many systems are compared on one
collection~\cite{carterette2012multiple}, and Boytsov, Belova and Westfall
worked through the practical choice of an adjustment~\cite{boytsov2013deciding};
the family-wise and false-discovery procedures they draw
on~\cite{holm1979simple,benjamini1995controlling} are the ones applied here.
Sakai has argued repeatedly for reporting effect sizes and confidence intervals
rather than significance alone~\cite{sakai2014statistical}, and Fuhr catalogued
the recurring errors that motivate that
argument~\cite{fuhr2018some}. Comparable arguments have been made for
natural-language evaluation: Dror and colleagues surveyed the choice of test for
language tasks~\cite{dror2018hitchhiker}, Card and colleagues showed how little
power common benchmark comparisons actually have~\cite{card2020little}, Reimers
and Gurevych showed that single-run scores conceal seed
variance~\cite{reimers2017reporting}, and Bouthillier and colleagues quantified
the variance sources a benchmark comparison must survive to mean
anything~\cite{bouthillier2021accounting}. The broader statistical-practice
critique this work sits inside is well
established~\cite{ioannidis2005why,simmons2011false,button2013power,wasserstein2016asa,gelman2014statistical,amrhein2019scientists}.

What this literature supplies is diagnosis and remedy at the level of the test:
choose a better test, correct for multiplicity, report an interval, collect more
topics. What remains to be supplied is a statement, in the units the evaluation
was actually recorded in, of how many individual gradings a particular reported
conclusion depends on. That is the reporting gap this paper addresses, by an
application of an established fragility construction to a paired grading unit.
The count is complementary to the interval and the corrected p-value, and it is
read beside them.

\subsection{Retrieval, question answering, and the evaluation those systems inherit}

The evaluation being re-read sits in the open-domain question-answering line
that treats a corpus as something a system retrieves from before it
answers~\cite{chen2017reading,rajpurkar2016squad,kwiatkowski2019nq,
joshi2017triviaqa,yang2018hotpotqa}. Sparse lexical matching remains the
baseline retriever~\cite{robertson2009bm25}; dense and late-interaction
retrievers, then retrieval-augmented generators, are the systems that replaced
it on those
benchmarks~\cite{lee2019latent,karpukhin2020dpr,khattab2020colbert,
guu2020realm,izacard2021fid,lewis2020retrieval,izacard2022atlas,
thakur2021beir}. Graph-structured retrieval is one more way of packaging the
same evidence~\cite{edge2024local}. What those papers report is almost always an
accuracy, or a ranking metric, on a question set whose items were graded once.
Assessor disagreement is an old fact of that
literature~\cite{voorhees2000variations,bailey2008relevance,artstein2008inter};
it is rarely converted into a count of how many of those grades a published
contrast actually depends on. The present paper asks that prior question of one
already-recorded comparison, and takes the retrievers as the records left them.

The same literature that diagnoses under-powered benchmark comparisons also
says how a small evaluation should be
read~\cite{cohen1992power,button2013power,card2020little,munafo2017manifesto,
head2015phacking,openscience2015repro,goodman2016mean,pineau2021improving,
wilkinson2016fair}. A paired grading count is one way of making that reading
concrete.

\subsection{Graded answers as the unit of evidence}

Treating an individual grading as the unit that can move also connects to work
on annotation reliability. Inter-coder agreement has long been understood as a
precondition for interpreting an evaluation rather than a formality to be
reported once~\cite{artstein2008inter}, and Voorhees's assessor study is an early
demonstration that the underlying judgments are genuinely
unstable~\cite{voorhees2000variations}. More recently, Northcutt, Athalye and
Mueller showed that label errors are pervasive enough in widely used test sets
to reorder published model rankings~\cite{northcutt2021pervasive}, and the
benchmark-critique literature has argued that scores are routinely read as more
meaningful than their construction
supports~\cite{bowman2021will,raji2021everything}. The count reported here makes
that connection quantitative for one evaluation: when the number of gradings a
conclusion depends on is smaller than the number of gradings that defensible
grading rules disagree about, the conclusion and the disagreement are the same
size, and the comparison between them is arithmetic rather than rhetorical.

\section{Methods}
%% Independent methods file. Not woven into main.tex prose.
%% Every quantity described here resolves to a hash-bound entry in the manifest
%% before it may be cited in Results.

\subsection{What was already on disk}

Nothing in this study was generated for it. No model was run, no answer was
produced, and no question was written. Every quantity below comes from records
written by runs that finished before this analysis began, and each record is
bound to the manuscript by its SHA-256 digest; the analysis code re-hashes each
file before reading it and stops the build on any mismatch.

The evaluation being re-read is small and simple. A question set of
\Questions{} questions was answered by a \Model{} instruction model at
temperature \Temperature{} under \ConditionsWord{} retrieval conditions --- no
retrieval, retrieval over a text-only knowledge graph, and retrieval over the
same graph with image-derived attributes attached --- and each answer was
graded right or wrong against a gold string. The latter two are the
graph-structured form of retrieval-augmented
generation~\cite{lewis2020retrieval,edge2024local} and differ only in whether
the graph carries attributes derived from images. Two such question sets exist, one
built around widely documented entities and one around less documented ones,
giving \Answers{} graded answers in total. Each set is split by a recorded
per-question flag into the \SubsetN{} questions whose answer depends on the
image and the \TextOnlyN{} that can be answered from text alone.

The design is paired: the same question is answered under every condition, so a
contrast between two conditions is decided entirely by the questions the two
conditions disagree on. Everything in this paper follows from that.

\subsection{The four counts a contrast reduces to}

For two conditions over a set of questions, each question falls into exactly one
of four cells: both conditions right, the first right alone, the second right
alone, or neither right. The difference in accuracy, the exact test and its
interval are all functions of those four counts and of nothing else.

Changing one grading --- deciding that one answer under one condition should
have been marked the other way --- moves one question from one cell to an
adjacent one. So the set of tables reachable from the observed one by $k$
grading changes is exactly the set reachable in $k$ steps of that move, and the
question this paper asks has an exact answer rather than an estimate.

\subsection{Two conclusions per contrast, and two distances}

A contrast can carry two different claims, and they are unseated by different
things. The first is a direction: which condition is ahead. The second is a
verdict: whether an exact test calls the difference significant at the
conventional threshold of \AlphaLevel{}. This study measures the distance to
each separately, because its central observation is that they are frequently not
the same distance.

The direction is unseated when the difference reaches zero or reverses. The
verdict is unseated when the p-value crosses the threshold, in whichever
direction it currently sits: a significant result is unseated by becoming
non-significant, and a non-significant one by becoming significant. The second
half of that is not symmetry for its own sake. A verdict of ``no effect'' is
reported and acted on exactly as a verdict of ``effect'' is, so it is entitled
to the same audit.

\subsection{Finding the minimum, rather than bounding it}

The search is exhaustive over reachable states rather than greedy over answers.
Because a contrast is fully described by four non-negative counts summing to
$n$, the reachable tables form a small graph, and a breadth-first walk outward
from the observed table returns the first state satisfying the stopping
condition. That state is at the true minimum distance, so the numbers reported
here are minima and not upper bounds on minima.

Where the minimum is one, the paper enumerates every individual grading that
would suffice on its own, so a reader can see which answers carry a conclusion
rather than being told how many. Several candidates do not make a conclusion
sturdier: it means there are several answers, any one of which carries it alone.

\subsection{What the test could have returned}

An exact test on paired binary data is a binomial test at one half on the
discordant pairs alone~\cite{mcnemar1947note}, which is the standard choice for
comparing two systems on a shared test set~\cite{dror2018hitchhiker}, so the
smallest p-value it can
return is fixed by the number of discordant pairs before any answer is graded.
That floor is reported for every contrast. A contrast whose floor lies above the
threshold could not have been called significant however the answers came out,
and its non-significance is therefore a fact about the question set rather than
about the conditions being compared.

Intervals are Clopper--Pearson on the discordant split mapped onto the paired
difference~\cite{clopper1934use}, matching the construction the records already
carry. Where a contrast has no discordant pairs the test is undefined and its
p-value is reported as one rather than as an error.

\subsection{Deletion, as a separate question}

Beside each flip distance the paper reports a leave-one-out
count~\cite{efron1979bootstrap}: how many of the $n$ single-question deletions
leave the direction, and separately the verdict, as it was. This measures
something different. A flip asks what a different grading of the same question
set would have done; a deletion asks what a different question set would have
done. A conclusion can be robust to one and fragile to the other, and both are
reported for every contrast rather than one standing in for the other.

\subsection{Multiplicity}

The records report \Contrasts{} paired contrasts and state that their p-values
are unadjusted. This paper adds the adjustment the record names as the one that
would raise the bar~\cite{holm1979simple} and reports it beside each raw value.
The adjustment is applied to the whole family of \Contrasts{}, because that is
the family that was computed; no subset of it is treated as though it had been
specified in advance.

\subsection{Checks the build enforces}

Several sentences in this paper are enforced as build-time checks rather than
asserted in prose, and the analysis code stops rather than reporting a number
that would contradict them.

Every one of the \Contrasts{} contrasts is recomputed from the per-answer rows
and checked against the recorded version on three quantities separately: the
number of questions, the difference in accuracy, and the exact p-value. The
discordant split is checked as well as the difference, because a table matching
on the difference alone could produce the same p-value by coincidence. The
recorded verdict is checked to follow from the recorded test rather than being
taken on its word.

The reshaping from per-answer rows into paired tables refuses a question that
was not answered under every condition, refuses a condition that answers the
same question twice, and refuses a question whose subset flag differs between
conditions --- each of which would let a pair be formed across the wrong things.
Each run's own summary is checked to still restate the rows in its own
directory, and the fixed measurement table the project recorded as its result is
checked to still come from those rows.

The regrade record is checked to contain exactly one contested grading, and to
be the one the manuscript names. The value the project recorded as not to be
used is checked to still be recorded that way. The tiers reported here as absent
or incomplete are checked to still be in that state, the blocked component is
checked to still be blocked with no score substituted for it, and the retrieval
result is checked to still carry the index state that keeps it a footnote.

\subsection{What is deliberately not computed}

No answer is regenerated and no grading is decided here; the contested grading
discussed below is one the records already contain, applied as the record states
it, not a judgement made by this paper. The reference benchmark this line of
work would eventually be measured against was never downloaded, the tier that
would hold that comparison is recorded as \PrimaryTier{} and the tier above it
as \SotaTier{}, and no value is imputed for either. The \BlockedComponent{}
component is recorded as blocked on a service that was unreachable, and its
scores \ForgedState{} invented in its place. Nothing in this paper substitutes
for any of them.


\begin{figure}[tbp]
\centering
\includegraphics[width=\linewidth]{fig1_plane.pdf}
\caption{\CapPlane}
\label{fig:plane}
\end{figure}

\section{Results}

Fig.~\ref{fig:plane} draws the headline conclusions in the plane the method
lives in, and it is worth reading before any number is quoted. The two right
panels are the same contrast, the same four counts and the same observed point,
read as two different claims; the outline around that point reaches \HelpK{}
gradings in one and \HelpVerdictK{} in the other. That difference belongs to
the two regions being measured to rather than to the data, and it is the whole
reason this paper reports two distances per contrast instead of one.

Table~\ref{tab:contrasts} gives all \Contrasts{} contrasts,
Table~\ref{tab:cells} gives the four counts each of them reduces to, and
Fig.~\ref{fig:flips} plots the two distances for each. Every recorded
difference and every recorded p-value reproduced from the per-answer rows, and
so did every discordant split; the build would have stopped otherwise, so what
follows is a re-reading of the project's own numbers. Both distances were
additionally recomputed from Propositions~\ref{prop:direction}
and~\ref{prop:verdict} and agreed with the search on every contrast.

\begin{figure}[tbp]
\centering
\includegraphics[width=\linewidth]{fig2_flips.pdf}
\caption{\CapFlips}
\label{fig:flips}
\end{figure}

\begin{table}[tbp]
\centering
\caption{\CapContrastsTable}
\label{tab:contrasts}
\scriptsize
\setlength{\tabcolsep}{4.5pt}
% Generated by the figure script from hash-bound evidence. Do not hand-edit.
\begin{tabular}{llrrrrrrrr}
\toprule
& & & & & & & \multicolumn{2}{c}{Gradings to unseat} & \\
\cmidrule(lr){8-9}
Questions & Contrast & $n$ & Disc. & Diff. (pt) & $p$ & Floor & Dir. & Verd. & Holm \\
\midrule
Famous, image & multimodal vs text & 18 & 7 & $-5.6$ & 1.000 & 0.016 & 1 & 5 & 1.000 \\
Famous, image & text vs none & 18 & 2 & $+0.0$ & 1.000 & 0.500 & 0 & 6 & 1.000 \\
Famous, image & multimodal vs none & 18 & 7 & $-5.6$ & 1.000 & 0.016 & 1 & 5 & 1.000 \\
Famous, text & multimodal vs text & 20 & 0 & $+0.0$ & 1.000 & 1.000 & 0 & 6 & 1.000 \\
Famous, text & text vs none & 20 & 4 & $+20.0$ & 0.125 & 0.125 & 4 & 2 & 1.000 \\
Famous, text & multimodal vs none & 20 & 4 & $+20.0$ & 0.125 & 0.125 & 4 & 2 & 1.000 \\
Famous, all & multimodal vs text & 38 & 7 & $-2.6$ & 1.000 & 0.016 & 1 & 5 & 1.000 \\
Famous, all & text vs none & 38 & 6 & $+10.5$ & 0.219 & 0.031 & 4 & 2 & 1.000 \\
Famous, all & multimodal vs none & 38 & 11 & $+7.9$ & 0.549 & $9.8 \times 10^{-4}$ & 3 & 4 & 1.000 \\
Long-tail, image & multimodal vs text & 18 & 8 & $+33.3$ & 0.070 & 0.008 & 6 & 1 & 0.984 \\
Long-tail, image & text vs none & 18 & 4 & $-11.1$ & 0.625 & 0.125 & 2 & 4 & 1.000 \\
Long-tail, image & multimodal vs none & 18 & 8 & $+22.2$ & 0.289 & 0.008 & 4 & 2 & 1.000 \\
Long-tail, text & multimodal vs text & 20 & 0 & $+0.0$ & 1.000 & 1.000 & 0 & 6 & 1.000 \\
Long-tail, text & text vs none & 20 & 12 & $+60.0$ & $4.9 \times 10^{-4}$ & $4.9 \times 10^{-4}$ & 12 & 4 & 0.008 \\
Long-tail, text & multimodal vs none & 20 & 12 & $+60.0$ & $4.9 \times 10^{-4}$ & $4.9 \times 10^{-4}$ & 12 & 4 & 0.008 \\
Long-tail, all & multimodal vs text & 38 & 8 & $+15.8$ & 0.070 & 0.008 & 6 & 1 & 0.984 \\
Long-tail, all & text vs none & 38 & 16 & $+26.3$ & 0.021 & $3.1 \times 10^{-5}$ & 10 & 2 & 0.319 \\
Long-tail, all & multimodal vs none & 38 & 20 & $+42.1$ & $4.0 \times 10^{-4}$ & $1.9 \times 10^{-6}$ & 16 & 6 & 0.007 \\
\bottomrule
\end{tabular}

\end{table}

\begin{table}[tbp]
\centering
\caption{\CapCellsTable}
\label{tab:cells}
\small
% Generated by the figure script from hash-bound evidence. Do not hand-edit.
\begin{tabular}{llrrrrrr}
\toprule
& & & \multicolumn{4}{c}{Questions answered correctly by} & \\
\cmidrule(lr){4-7}
Questions & Contrast & $n$ & Both & Arm & Ref. & Neither & Disc. \\
\midrule
Famous, image & multimodal vs text & 18 & 4 & 3 & 4 & 7 & 7 \\
Famous, image & text vs none & 18 & 7 & 1 & 1 & 9 & 2 \\
Famous, image & multimodal vs none & 18 & 4 & 3 & 4 & 7 & 7 \\
Famous, text & multimodal vs text & 20 & 20 & 0 & 0 & 0 & 0 \\
Famous, text & text vs none & 20 & 16 & 4 & 0 & 0 & 4 \\
Famous, text & multimodal vs none & 20 & 16 & 4 & 0 & 0 & 4 \\
Famous, all & multimodal vs text & 38 & 24 & 3 & 4 & 7 & 7 \\
Famous, all & text vs none & 38 & 23 & 5 & 1 & 9 & 6 \\
Famous, all & multimodal vs none & 38 & 20 & 7 & 4 & 7 & 11 \\
Long-tail, image & multimodal vs text & 18 & 3 & 7 & 1 & 7 & 8 \\
Long-tail, image & text vs none & 18 & 3 & 1 & 3 & 11 & 4 \\
Long-tail, image & multimodal vs none & 18 & 4 & 6 & 2 & 6 & 8 \\
Long-tail, text & multimodal vs text & 20 & 20 & 0 & 0 & 0 & 0 \\
Long-tail, text & text vs none & 20 & 8 & 12 & 0 & 0 & 12 \\
Long-tail, text & multimodal vs none & 20 & 8 & 12 & 0 & 0 & 12 \\
Long-tail, all & multimodal vs text & 38 & 23 & 7 & 1 & 7 & 8 \\
Long-tail, all & text vs none & 38 & 11 & 13 & 3 & 11 & 16 \\
Long-tail, all & multimodal vs none & 38 & 12 & 18 & 2 & 6 & 20 \\
\bottomrule
\end{tabular}

\end{table}

\subsection{The three prespecified conclusions}

The three prespecified conclusions are in Table~\ref{tab:prereg}, and they are
taken in turn.

\begin{table}[tbp]
\centering
\caption{\CapPreregTable}
\label{tab:prereg}
\small
% Generated by the figure script from hash-bound evidence. Do not hand-edit.
\begin{tabular}{lrrrr}
\toprule
 & As & Gradings & Single-grading & Deletions \\
Conclusion & measured & to unseat & candidates & survived \\
\midrule
Multimodal is behind text, famous & $-5.6$ & 1 & 19 & 14 of 18 \\
Multimodal is ahead of text, long-tail & $+33.3$ & 6 & --- & 18 of 18 \\
Text beats no retrieval, long-tail & $+60.0$, $4.9 \times 10^{-4}$ & 4 & --- & 20 of 20 \\
\midrule
Same contrast, verdict not direction & 0.070 & 1 & 12 & 17 of 18 \\
\bottomrule
\end{tabular}

\end{table}

The first is that on the image-dependent questions of the widely-documented set,
the multimodal condition is behind the text-only one: \HarmMM{} against
\HarmText{} over \SubsetN{} questions, a difference of \HarmDiff{} percentage
points on \HarmDiscordant{} discordant pairs, with an exact p-value of
\HarmP{}. The minimum flip set is \HarmK{} grading, and the enumeration is
worth more than the count: \HarmCandidates{} of the \HarmGradings{} individual
gradings on that subset would each, alone, be enough to remove the difference or
reverse it. Each of those \HarmCandidates{} candidate gradings is therefore
sufficient to erase its reported direction.

The second is that on the image-dependent questions of the less-documented set,
the multimodal condition is ahead: \HelpMM{} against \HelpText{}, a difference
of \HelpDiff{} percentage points, exact p-value \HelpP{}, which the records
report as not significant and which is read here as a recorded difference
rather than as an established effect. Its direction needs \HelpK{} gradings to
reverse. That is the sturdiest direction among the three at this subset size,
and it is still \HelpK{} answers out of \SubsetNWord{}.

The third survives Holm correction across the family of \Contrasts{} under the
stored grading rule: on the text-answerable questions of the less-documented
set, the text-retrieval condition beats no retrieval by \SigDiff{} percentage
points with an exact p-value of \SigP{}, or \SigHolm{} after correction. Under
the numeric-tolerant nested rule its adjusted p-value is \NumericHolmLT{} and
it does not survive; survival is a property of the grader and the family, not
of the contrast alone. Its verdict needs \SigK{} gradings and its direction
needs \SigSignK{}. This is what a comparatively secure result looks like at
this scale, and its content is modest: it says that supplying a fact makes a
model more likely to state that fact.

The first conclusion is also the cleanest illustration of
Proposition~\ref{prop:direction}. Its \HarmDiff{}-point difference on \SubsetN{}
questions is a lead of one question in the discordant cells, and the proposition
says the distance to its reversal is that lead. A reader shown the percentage
and the sample size has been shown the count already, in a form that requires
division to recover; a reader shown the count needs nothing.

The flip ladder of Fig.~\ref{fig:flips} counts how many individual gradings
could alter a conclusion; the same contrasts can also be read on their accuracy
scale. Fig.~\ref{fig:effect-forest}
puts the recorded difference and its exact paired interval beside one another
for all \Contrasts{} rows. It is a descriptive re-encoding of the four-cell
tables rather than a second test: the compact key suppresses item text, and $m$
makes the number of informative questions visible next to the percentage-point
difference. The spread of intervals is therefore read together with the flip
counts, as a description of the recorded tables rather than as evidence for a
retrieval effect.

\begin{figure}[tbp]
\centering
\includegraphics[width=\linewidth]{fig8_effect_forest.pdf}
\caption{\CapEffectForest}
\label{fig:effect-forest}
\end{figure}

\subsection{The direction and the verdict are not equally secure}

The fourth row of Table~\ref{tab:prereg} is not one of the three, and it is the
result this paper would keep if it could keep only one. The
\HelpDiff{}-point difference (exact \(p=\HelpP{}\), reported as not
significant) whose direction needs \HelpK{} gradings has a
verdict --- not significant at \AlphaLevel{} --- that needs \HelpVerdictK{}, and
\HelpVerdictCandidates{} distinct single gradings would each be enough to make
it significant. Its floor is \HelpFloor{}, so it was capable of significance and
simply did not reach it.

The same four counts, read two ways, give a claim that \HelpK{} answers could
not shake and a claim that \HelpVerdictK{} answer could. Which of the two a
reader takes away is decided by the conventions of the table it is printed in,
and a table reporting $p = \HelpP{}$ with no further annotation invites the
reading that has the weaker support.

Across the study this is the common case rather than the exception.
Fig.~\ref{fig:decouple} plots the two distances against each other for the
whole family. In \DecoupledCount{} of the \Directional{} contrasts that have a
direction to reverse, the two differ by \DecoupledGap{} gradings or more, and
they differ in both directions: \FragileSign{} contrasts have a direction that
\FragileMax{} gradings or fewer would reverse, \FragileVerdict{} have a verdict
that \FragileMax{} or fewer would flip, and no contrast is in both groups.
Points fall on both sides of the diagonal. There is no general rule that one is
sturdier. There is only the fact that they are different questions, and that
reporting one number for both loses the distinction.

\begin{figure}[tbp]
\centering
\includegraphics[width=0.51\linewidth]{fig3_decouple.pdf}
\caption{\CapDecouple}
\label{fig:decouple}
\end{figure}

The reason the two can separate at all is structural rather than accidental, and
Fig.~\ref{fig:plane} shows it. The direction fails on a half-plane, so its
distance is the lead in discordant pairs and nothing else. The verdict fails on
a region whose boundary bends, because the exact test reads the discordant count
as well as the split, so a flip that moves a pair from one side to the other and
a flip that removes a pair altogether are worth different amounts to it. Two
distances measured to two differently shaped regions from the same point have no
reason to agree, and the finding here is that at this scale they routinely do
not.

\subsection{The \RegradeContested{} answers on which the nested rules disagree}

Fig.~\ref{fig:regrade} makes the first conclusion concrete. The same
\RegradeAnswers{} answers were scored by \Graders{} nested string rules, which
agree on all but \RegradeContested{}. Table~\ref{tab:graders} gives the
resulting aggregate accuracies: \RegradeIdenticalRows{} of its twelve rows are
identical across the three rules, and \RegradeMovingRows{} move. One of those
moving rows is the multimodal condition on the image-dependent split of the
widely-documented set, and it is the row Fig.~\ref{fig:regrade} isolates. The
other is the no-retrieval condition on the less-documented text split. Question
keys, answer strings, and gold values remain in the evidence rather than in the
manuscript. The image-subset disagreement is one of the \HarmCandidates{}
gradings the search identified.

\begin{figure}[tbp]
\centering
\includegraphics[width=0.51\linewidth]{fig4_regrade.pdf}
\caption{\CapRegrade}
\label{fig:regrade}
\end{figure}

\begin{table}[tbp]
\centering
\caption{\CapGradersTable}
\label{tab:graders}
\small
% Generated by the figure script from hash-bound evidence. Do not hand-edit.
\begin{tabular}{lllrrr}
\toprule
& & & \multicolumn{3}{c}{Accuracy under each retrieval condition} \\
\cmidrule(lr){4-6}
Questions & Split & Grader & none & text & multimodal \\
\midrule
Famous & image & strict match & 0.444 & 0.444 & 0.389 \\
Famous & image & substring match & 0.444 & 0.444 & 0.389 \\
Famous & image & numeric tolerance & 0.444 & 0.444 & 0.444 \\
Famous & text & strict match & 0.800 & 1.000 & 1.000 \\
Famous & text & substring match & 0.800 & 1.000 & 1.000 \\
Famous & text & numeric tolerance & 0.800 & 1.000 & 1.000 \\
Long-tail & image & strict match & 0.333 & 0.222 & 0.556 \\
Long-tail & image & substring match & 0.333 & 0.222 & 0.556 \\
Long-tail & image & numeric tolerance & 0.333 & 0.222 & 0.556 \\
Long-tail & text & strict match & 0.400 & 1.000 & 1.000 \\
Long-tail & text & substring match & 0.400 & 1.000 & 1.000 \\
Long-tail & text & numeric tolerance & 0.550 & 1.000 & 1.000 \\
\bottomrule
\end{tabular}

\end{table}

Applying it as the record states moves the multimodal condition from
\HarmMM{} to \RegradedMM{} on those \SubsetN{} questions and the contrast from
\HarmDiff{} to \RegradedDiff{} percentage points. The reported direction does
not weaken; it disappears. Under one rule the multimodal condition is behind,
under the other the two are exactly level, and the two rules differ on how to
read one numeric string. The choice between them is left open here.

The three rules are nested string matches rather than independent graders:
exact token, substring, then a numeric window that also fires on year answers.
Two local instruction-model judges, scoring the stored strings and nothing
else, agree with the stored rule at Cohen's \(\kappa=\LLMKappaQwen{}\) and
\(\kappa=\LLMKappaLlama{}\), and with each other at \(\kappa=\LLMKappaBoth{}\).
That agreement is a consistency check on the stored strings rather than an
independent rating of the answers: the judges read the same strings the rules
read, and what they confirm is that the stored rule's verdicts are the ones a
reader of those strings would also reach, which is a different thing from its
being the right rule.

The choice between the two graders is left open because it is a real
methodological question rather than a formality: whether a tolerance is
appropriate depends on what the question was asking, and that is a judgement
about the item. The point of this paper is upstream of the choice. A reported
direction that is decided by an unstated grading convention was never a
property of the conditions being compared, and the way to notice that before
publishing is to count, because the count is what makes a grading convention
visible as a load-bearing part of the result.

That the flip turned up in an independent regrade rather than in the search is
also the reason to trust the search. The two procedures had nothing to do with
each other: one enumerated gradings that would change a conclusion, the other
re-scored answers under different string-matching rules, and they met at the
same answer.

Fig.~\ref{fig:flip-summary} collects the three pre-registered claims under
both graders. The stored strict grader gives direction distances
\SummaryHarmStrict{} and \SummaryHelpStrict{} for the harm and help contrasts
and a verdict distance of \SummarySigStrict{} for the long-tail text contrast.
Repeating the same per-contrast calculation under the numeric-tolerant grader
gives \SummaryHarmNumeric{}, \SummaryHelpNumeric{} and \SummarySigNumeric{},
respectively. These are raw per-contrast distances conditioned on the grader.
Holm survival is a separate family-level condition and is reported for each
grader on its own, so the panel conditions distances on the grader without
treating them as family-adjusted quantities.

\begin{figure}[tbp]
\centering
\includegraphics[width=0.51\linewidth]{fig10_flip_regrade_summary.pdf}
\caption{\CapFlipSummary}
\label{fig:flip-summary}
\end{figure}

\subsection{Counting candidate gradings without publishing row identities}

Where a conclusion's distance is one, the evidence enumerates every candidate
grading, and the manuscript reports the count. Table~\ref{tab:candidates} gives
that aggregate count for both such conclusions and marks the row identities as
withheld. This preserves the fragility result while keeping question keys,
answer strings, and gold values out of the manuscript. A reader can audit the
count from the redacted evidence release that accompanies it, and no row-level
material is required to interpret the figure or the conclusion.

\begin{table}[tbp]
\centering
\caption{\CapCandidatesTable}
\label{tab:candidates}
\small
% Generated by the figure script from hash-bound evidence. Do not hand-edit.
\begin{tabular}{lrp{0.52\linewidth}}
\toprule
Condition regraded & Count & Questions whose single regrading suffices \\
\midrule
\multicolumn{3}{l}{\emph{Multimodal is behind text, widely-documented set: the direction}} \\
\quad text & 8 & I2, I3, I4, I6, I8, I11, I13, I17 \\
\quad multimodal & 11 & I1, I2, I6, I7, I8, I12, I14, I15, I16, I17, I18 \\
\midrule
\multicolumn{3}{l}{\emph{Less-documented image-dependent difference: the verdict}} \\
\quad text & 4 & LI3, LI4, LI8, LI16 \\
\quad multimodal & 8 & LI2, LI6, LI7, LI8, LI13, LI14, LI15, LI17 \\
\bottomrule
\end{tabular}

\end{table}

The proportions are worth reading directly. The reported finding that the
multimodal condition trails the text-only one on the image-dependent questions
of the widely-documented set rests on \HarmGradings{} individual gradings, of
which \HarmCandidates{} would each erase it alone. The verdict on the opposing
contrast has \HelpVerdictCandidates{} such answers. Each count is a count of
adjudications rather than a confidence statement: it says how many distinct
gradings could carry the conclusion, and the identity of each stays in the
evidence.

\subsection{What each test could have returned, and what survives correction}

Fig.~\ref{fig:floor} reports a quantity that is available before any answer is
graded. An exact paired test is a binomial test on the discordant pairs
alone~\cite{mcnemar1947note}, so the smallest p-value it can produce is fixed by
how many questions the two conditions disagreed on. With three discordant pairs
the most extreme result attainable is a two-sided p-value of one quarter, and no
grading of those three answers can do better.

\begin{figure}[tbp]
\centering
\includegraphics[width=0.51\linewidth]{fig5_floor.pdf}
\caption{\CapFloor}
\label{fig:floor}
\end{figure}

Table~\ref{tab:design} makes that arithmetic a lookup. Its first three columns
depend on nothing about this evaluation: they say what any paired comparison at
this threshold can return given how many items the two systems disagreed on, and
they are the content of Corollary~\ref{cor:floor}. Below
\SignificanceNeedsM{} discordant pairs the column reads unreachable, because no
assignment of that many disagreements reaches \AlphaLevel{}. The last column is
this evaluation's entry in the table: \Incapable{} of its \Contrasts{}
contrasts sit in those rows.

\begin{table}[tbp]
\centering
\caption{\CapDesignTable}
\label{tab:design}
\small
% Generated by the figure script from hash-bound evidence. Do not hand-edit.
\begin{tabular}{rrlr}
\toprule
Discordant & Smallest attainable & Majority needed & Contrasts \\
pairs $m$ & $p$-value & for significance & with this $m$ \\
\midrule
0 & 1.000 & unreachable & 2 \\
1 & 1.000 & unreachable & 0 \\
2 & 0.500 & unreachable & 1 \\
3 & 0.250 & unreachable & 0 \\
4 & 0.125 & unreachable & 3 \\
5 & 0.062 & unreachable & 0 \\
6 & 0.031 & 6 of 6 & 1 \\
7 & 0.016 & 7 of 7 & 3 \\
8 & 0.008 & 8 of 8 & 3 \\
9 & 0.004 & 8 of 9 & 0 \\
10 & 0.002 & 9 of 10 & 0 \\
11 & $9.8 \times 10^{-4}$ & 10 of 11 & 1 \\
12 & $4.9 \times 10^{-4}$ & 10 of 12 & 2 \\
13 & $2.4 \times 10^{-4}$ & 11 of 13 & 0 \\
14 & $1.2 \times 10^{-4}$ & 12 of 14 & 0 \\
15 & $6.1 \times 10^{-5}$ & 12 of 15 & 0 \\
16 & $3.1 \times 10^{-5}$ & 13 of 16 & 1 \\
17 & $1.5 \times 10^{-5}$ & 13 of 17 & 0 \\
18 & $7.6 \times 10^{-6}$ & 14 of 18 & 0 \\
19 & $3.8 \times 10^{-6}$ & 15 of 19 & 0 \\
20 & $1.9 \times 10^{-6}$ & 15 of 20 & 1 \\
\bottomrule
\end{tabular}

\end{table}

Those \Incapable{} contrasts were reported as non-significant, and they could
not have been reported as anything else. For them the phrase carries no
information about the conditions being compared: it is a restatement of how
rarely the two conditions disagreed, which is a property of the question set.

A second reading of the same figure is the \AtFloor{} contrasts sitting on the
diagonal, where every discordant pair fell the same way and the test returned
the most extreme value available to it. Of those, \AtFloorNS{} are
non-significant regardless. A contrast in which one condition won every single
question the two conditions disagreed about, and which is nonetheless reported
as no evidence of a difference, is the clearest available demonstration that a
threshold verdict at this scale is a statement about $n$ before it is a
statement about the conditions.

Of the \Contrasts{} contrasts, \SignificantRaw{} reach the threshold unadjusted
and \SurvivesHolm{} survive correction across the family~\cite{holm1979simple}
under the stored rule. The records state that their p-values are unadjusted and
that a correction would raise the bar; this paper applies that correction. The
survivors are all differences against no retrieval at all, and none of them is
a comparison between the two retrieval conditions, which is the comparison the
evaluation was built to make. Two of those \SurvivesHolm{} rows are the same
long-tail text table counted twice. Under the numeric-tolerant nested rule of
the same regrade record that table moves from Holm-adjusted \(p=\SigHolm{}\) to
\(p=\NumericHolmLT{}\) and does not survive; \NumericHolmSurvivors{} contrast in
the family of \Contrasts{} then remains. Survival is therefore a joint fact
about the grading rule and the family definition rather than a property of the
evaluation alone.

\subsection{How often a significant verdict rests on one grading}

The closed form makes a sampling experiment cheap. Fig.~\ref{fig:operating}
reports \SimTables{} paired Bernoulli tables at the \(n\) of this evaluation and
at smaller and larger designs, \SimReps{} replicates per cell. Conditional on
\(p\le\AlphaLevel{}\), the share of tables whose verdict distance is one is
\FragileGivenSigNEighteen{} at \(n=\SubsetN{}\) and
\FragileGivenSigNThirtyEight{} at \(n=\Questions{}\); both fall toward zero as
\(n\) grows. Among significant tables at \(n=\Questions{}\), Spearman correlation
of the verdict distance with exact \(p\) is \SpearmanSigNThirtyEight{}, so the
count tracks \(p\) closely on that side of the threshold and still inverts
adjacent tables: \KappaVerInversionsLow{}--\KappaVerInversionsHigh{} inversions
among \KappaVerSigTables{} significant states. Under this model a significant
verdict at these sizes is typically a one-grading result unless the effect is
large. That is an operating characteristic of the statistic --- a description
of how the verdict distance behaves at small \(n\) --- and it is read as such,
separately from the findings about the evaluation.

\begin{figure}[tbp]
\centering
\includegraphics[width=0.51\linewidth]{fig9_operating.pdf}
\caption{\CapOperating}
\label{fig:operating}
\end{figure}

\subsection{The same family read against every threshold}

Everything above is read against \AlphaLevel{}, which is the level the records
used, and a reader is entitled to ask how much of it is a fact about the
evaluation and how much a fact about that number. Fig.~\ref{fig:alpha}
recomputes the whole family at \SweepLevels{} thresholds between \SweepLow{} and
\SweepHigh{}.

\begin{figure}[tbp]
\centering
\includegraphics[width=0.51\linewidth]{fig6_alpha.pdf}
\caption{\CapAlpha}
\label{fig:alpha}
\end{figure}

The counts move as they must: a looser threshold admits more contrasts and
leaves fewer incapable of reaching it. What does not move is the shape. Across
that whole range the number of contrasts that could not have reached the
threshold whatever the answers were never falls below \IncapableSweepMin{},
because that number is fixed by how rarely the two conditions disagreed rather
than by where the line is drawn, and the median number of gradings that would
move a verdict stays between \MedianKSweepMin{} and \MedianKSweepMax{}.
Choosing a different conventional threshold would change which contrasts are
called significant. It would not change the observation that the call is being
made on a handful of individual adjudications.

The sweep is cheap because the closed form is exact. Recomputing every distance
at \SweepLevels{} thresholds is one pass per contrast per level. The search
remains the build's unit test at the recorded threshold, where it is checked
against the closed form on every contrast.

\subsection{Deletion asks a different question from re-grading}

Fig.~\ref{fig:loo} reports the leave-one-out count beside the flip distance,
and the two disagree in an instructive way. Most contrasts survive every single
deletion: the direction and the verdict are unchanged in all $n$ recomputations.
Deletion is a blunt instrument here, because removing a concordant question
changes neither the discordant split nor the sign, and most questions are
concordant.

\begin{figure}[tbp]
\centering
\includegraphics[width=0.51\linewidth]{fig7_loo.pdf}
\caption{\CapLeaveOneOut}
\label{fig:loo}
\end{figure}

The exceptions are the contrasts already identified. The direction of the first
conclusion survives \HarmLooSign{} of \SubsetN{} deletions, so \HarmLooBreaks{}
individual questions each carry it by their presence as well as by their
grading. The verdict on the second survives \HelpLooVerdict{} of \SubsetN{}.

The two measurements answer different counterfactuals and should be read as
such. A flip asks what a different adjudication of the same evaluation would
have concluded, which is a question about grading policy and about annotator
disagreement~\cite{northcutt2021pervasive}. A deletion asks what a slightly
different question set would have concluded, which is a question about sampling.
An evaluation can be exposed on one and not the other, and reporting only the
one that happens to look better is a degree of freedom worth closing in
advance~\cite{simmons2011false}.

\section{Discussion}

The count is old and the target is new. Counting how many outcomes would have
to change to overturn a result is the fragility
index~\cite{feinstein1990unit,walsh2014statistical}, introduced for randomised
trials, where it revealed that a large share of results in the clinical
literature turn on a handful of events --- one mechanism behind the more general
observation that small studies produce unreliable
findings~\cite{ioannidis2005why}. The construction transfers to machine
learning evaluation almost unchanged, and the reasons to want it are stronger
here: an evaluation item is graded by a rule that someone wrote, benchmark labels
are known to carry errors at rates that change published
rankings~\cite{northcutt2021pervasive}, and the evaluation sets that appear in
practice are frequently small enough that the count is in single
digits~\cite{card2020little}. The main criticism of the index in its original
setting is that it is a monotone function of the p-value and so adds
nothing~\cite{carter2017fragility}. That criticism has less force here, for a
reason this study can now count: among \KappaVerSigTables{} significant tables
at \(n=\Questions{}\), ordering by exact \(p\) produces
\KappaVerInversionsLow{}--\KappaVerInversionsHigh{} adjacent inversions of the
verdict distance, depending on how ties in \(p\) are broken. The count is being
applied to the direction of a difference as well as to the verdict on it, the
two are not monotone in each other, and \DecoupledCount{} of the \Directional{}
contrasts that have a direction to reverse separate them by \DecoupledGap{}
gradings or more. The construction itself is Feinstein's; what this paper adds
is the paired grading unit it is counted in and the second claim, the
direction, that it is measured against.

A contrast carries two conclusions, and only one of them is usually audited.
Significance is routinely questioned and directions are routinely believed. This
study found the opposite pattern at least as often. The clearest case is the
contrast whose direction \HelpK{} gradings could not reverse and whose verdict
\HelpVerdictK{} grading would flip, but the general point is that these are
separate claims with separate supports, and a results table reports them in a
single row as though they stood or fell together. The recommendation implied is
small: print the two counts. They cost nothing to compute, they are exact, and
they say what a conditional interval says while remaining in the units the
reader can inspect --- answers.

Non-significance is a finding only where significance was reachable. Of the
\Contrasts{} contrasts here, \Incapable{} could not have been significant at
any grading of their answers. Reporting them as null results looks conservative
and is the opposite: it attributes to the conditions a property of the question
set. The floor is one line of arithmetic from the discordant count and belongs
beside any negative result from a paired design, in the same way that a power
calculation belongs beside a negative result from an unpaired
one~\cite{card2020little,amrhein2019scientists}. The stronger version of the
point is the contrast in which every disagreement fell one way and the verdict
was still negative: unanimity among the informative questions was not enough,
because there were \AtFloorNSDiscordant{} of them.

A grading convention is a researcher degree of freedom. String matching against
gold answers looks like a mechanical step and is a choice among rules. Strict
matching, substring matching and numeric tolerance are \Graders{} defensible
rules; they disagree about \RegradeContested{} answers here,
\RegradeImageFlip{} of which sits on the image-subset contrast and decides its
reported direction, while the rest decide whether the long-tail text pair
survives correction across the family. The literature on undisclosed analytic
flexibility~\cite{simmons2011false,gelman2014statistical} is usually read as
being about model selection and outlier handling; grading rules belong on the
same list, and in small evaluations they may dominate it. The regrade in this
study is the practice that made the point visible, and it cost one pass over the
stored answers with no model in the loop.

The concrete proposal is one column, reported with the claim. Beside each
reported contrast, print the smallest number of gradings that would overturn
its direction and the smallest that would move its verdict. Both are
exhaustively computable at these sizes. A reader seeing \HarmK{} in that column
knows immediately what kind of object they are being shown, without needing to
reconstruct the discordant table from an accuracy and a sample size --- which is
what it currently takes, and which is why the question is rarely asked. This is
the same instinct behind reporting effect sizes and intervals rather than
verdicts~\cite{wasserstein2016asa,amrhein2019scientists}, expressed in the unit
that a small evaluation is actually made of.

Counting gradings makes fragility visible; removing it is a separate task. An
evaluation whose conclusions rest on single answers needs more questions, not
better reporting of how few it had~\cite{card2020little,bowman2021will}. The
count is also silent on whether the question set measures what it claims to
measure, which at these sizes is the larger risk~\cite{raji2021everything}. Its
value is specific: it is cheap, it is exact, it applies to records that already
exist, and it prevents one particular error, publishing an arithmetic accident
in the register of a finding.

\section{Scope of the conclusions}

The conclusions of this paper are about how an evaluation is reported, and
three boundaries fix their scope.

The first concerns the retrieval conditions. Whether adding image-derived
attributes to retrieval helps or harms is a question these records do not
settle, and this paper leaves it open. The \HelpDiff{}-point contrast has an
exact p-value of \HelpP{} and is read here as a recorded difference rather than
as an established effect; the \HarmDiff{}-point contrast is \HarmK{} grading
from disappearing. Both enter as objects whose fragility is measured, and
neither is offered as evidence about the conditions compared.

The second concerns reference benchmarks, on which this paper reports no
performance. The benchmark this line of work would be measured against was
never obtained; that comparison is recorded as \PrimaryTier{}, the comparison
against the strongest published method that would follow it as \SotaTier{}, and
no value is imputed for either. The \BlockedComponent{} component is recorded
as blocked on a service that was unreachable, and its scores \ForgedState{}
invented in its place; the \RetrievalCorrect{} of \RetrievalN{} retrieval figure
is a footnote on a lookup over an index of \IndexFiles{} files rather than an
accuracy, and the value the records mark as not to be used appears nowhere
above. Each of these is a state the build checks on every run rather than an
assurance offered in prose.

The third concerns what a minimum flip set of one means. It means the
conclusion is supported by one grading, and it does not mean the conclusion is
false. A true effect measured on \SubsetNWord{} questions can rest on one
grading too; the count does not separate those cases, and what it establishes
is what separating them would take, which is more questions.

\section{Conclusion}

One completed evaluation of \Questions{} questions under \ConditionsWord{} retrieval
conditions was re-read by counting rather than by re-running. For each of its
\Contrasts{} paired contrasts, the smallest number of individual gradings that
would overturn the direction, and separately the verdict, was found by
exhaustive search over the reachable tables and confirmed by a closed form.

The three conclusions named in advance came out at \HarmK{}, \HelpK{} and
\SigK{} gradings, so the prespecified kill \KillState{}. The smallest,
\SmallestK{}, belongs to the reported finding that the multimodal condition
trails the text-only one: \HarmCandidates{} of the \HarmGradings{} gradings on
that subset would each erase it alone, and an independent regrade of the same
answers disagrees about \RegradeContested{} of them, \RegradeImageFlip{} of which moves the
harm difference from \HarmDiff{} to \RegradedDiff{} percentage points. The
direction of the opposing
\HelpDiff{}-point difference (exact \(p=\HelpP{}\), reported as not significant)
needs \HelpK{} gradings while the verdict attached to it needs \HelpVerdictK{},
and across the study those two distances differ by
\DecoupledGap{} or more in \DecoupledCount{} of the \Directional{} contrasts
that have a direction at all. No grading of the answers could have made
\Incapable{} of the contrasts significant, and survival of Holm correction is a
joint fact about the grading rule and the family: \SurvivesHolm{} contrasts
survive under the stored rule and \NumericHolmSurvivors{} under the
numeric-tolerant one.

The recommendation is one column in a results table: beside each contrast, the
number of gradings that would overturn its direction and the number that would
move its verdict. It is exactly computable, it costs nothing at these sizes, and
it converts a question a reader cannot currently ask into a number they can
read. When the count is large, nothing has been lost by printing it. When it is
\HarmK{}, what has been learned is that the row is reporting an adjudication of
one answer, and that is worth knowing before the table is published rather than
after.

\section*{Funding}

%% Funding and competing-interest wording adopted from the owner's packaging
%% script (scripts/build_springer_submission_package.py), shipped in every
%% packet since 2026-08-28 (ruling R16, 2026-09-09). FUNDING_AND_CONFLICTS
%% stays OPEN for the owner's final read. The comment sits after the first
%% \section* so the Discussion slice used by the checks is unchanged.
No external funding was received for this study.

\section*{Competing interests}

The author declares that they have no competing financial or non-financial
interests that are directly or indirectly related to the work submitted for
publication.

%% Draft-only (ruling R15, 2026-09-09): prints only when \DRAFTTODO is
%% defined; the owner supplies the sentence or deletes the heading. Kept below
%% the first \section* so the Discussion slice used by the checks is unchanged.
\ifdefined\DRAFTTODO
\section*{Acknowledgements}

\ownertodo{acknowledgements of people, institutions and computing resources
are to be supplied by the author; no acknowledgement is recorded in the
project ledger.}
\fi

\section*{Ethics approval and consent}

This study is a secondary analysis of numeric records produced by earlier
computational runs. It involved no human participants and no animals and
processed no personal data, so ethics approval, consent to participate and
consent for publication do not apply. No new evaluation answers were generated;
two local instruction-model judges scored the stored answer strings as a
consistency check and wrote no new answers. The questions concern publicly
documented landmarks and comparable entities. Grading disagreements are
reported in aggregate, and the manuscript prints no question keys, answer
strings or gold values. One consideration bears directly on the subject matter:
the paper's central observation is that a reported difference between two
systems can be decided by a single grading decision, and the same is true of
evaluations used to support claims about deployed systems, where a fragile
result may be read as a safety or capability assurance. Reporting the count is
offered partly for that reason.

\section*{Data availability}

The per-answer records for both question sets, the summaries each run wrote for
itself, the recorded paired contrasts with their exact tests and intervals, the
independent regrade under three graders, the fixed measurement table, and the
recorded status of each planned comparison, including those that are absent or
blocked, accompany the manuscript source. Each artifact is recorded under an
evidence ID in a manifest, and the analysis code verifies the bound bytes
before reading them, so the build stops rather than silently reporting stale
numbers. The knowledge graph and question sets are small local artifacts
distributed with the records; no third-party benchmark is redistributed, and
the model used to produce the answers is a third-party release identified by
name rather than copied.
%% Generated by scripts/release_targets.py from the release ledger.
%% Do not edit: an identifier typed here would not be checked against
%% anything, which is exactly the failure the ledger exists to prevent.
%% Ledger state: github=CREATED, zenodo=DEPOSITED
That material is deposited under the persistent identifier \href{https://doi.org/10.5281/zenodo.22647028}{10.5281/zenodo.22647028}, and the development repository is at \url{https://github.com/PeterPonyu/minimum-flip-set-audit}.


\section*{Code availability}

The analysis code that produces every number, table and figure in this
manuscript accompanies the manuscript source in the same deposit and the same
repository as the records, together with the exhaustive-search unit test that
checks the closed form on every reachable table and the build script that
verifies the bound evidence before reading it.

\section*{Author contributions}

\ifdefined\ANONYMOUS
The sole author performed conceptualization, methodology, software, validation,
formal analysis, investigation, visualization, writing---original draft, and
writing---review and editing.
\else
Zeyu Fu: conceptualization, methodology, software, validation, formal analysis,
investigation, visualization, writing---original draft, and writing---review
and editing.
\fi

\section*{Declaration of generative AI use}

Generative AI tools assisted with code preparation and language editing. Every
number printed above is emitted programmatically from artifacts tied to the
manuscript by evidence binding rather than transcribed into prose. Structural
claims --- including the paired differences, discordant splits, exact tests,
regrade record, and the handling of blocked comparisons --- are enforced by
build-time checks that stop the build. The evaluated responses were
pre-existing records and were not regenerated for this analysis. Responsibility
for the content rests with the author.
